% L4c lecture companion deck.
% Mirrors the lecture notebook: shortest-path formulation and relaxation;
% Dijkstra; Bellman-Ford; worked comparison; negative weights; lab bridge.
% The project is self-contained: styles, seal, diagrams, and build rules are local.
\documentclass[aspectratio=169,11pt]{beamer}
\usepackage{vnslides}
\usepackage{tabularx}

\renewcommand{\vnlecturecode}{L4c}
\newcolumntype{Y}{>{\raggedright\arraybackslash}X}

\tikzset{
  graphvertex/.style={draw=vnslideink,circle,fill=white,line width=0.7pt,
    minimum size=8mm,inner sep=0pt,font=\small\vnmedium},
  graphsource/.style={graphvertex,draw=vnaccent,line width=1.2pt},
  graphedge/.style={draw=vnslidebody,line width=0.8pt,-{Latex[length=2.5mm]}},
  graphoptimal/.style={draw=vnaccent,line width=1.4pt,-{Latex[length=2.7mm]}},
  graphnegative/.style={draw=vnslidelink,line width=1.2pt,-{Latex[length=2.7mm]}},
}

% Keep the worked graph visible beside each trace. Arguments are the five weights.
\newcommand{\tracegraph}[5]{%
  \begin{tikzpicture}[x=0.70cm,y=0.82cm]
    \node[graphsource] (v1) at (0,1.3) {1};
    \node[graphvertex] (v3) at (2,0) {3};
    \node[graphvertex] (v2) at (2,2.6) {2};
    \node[graphvertex] (v4) at (4.3,1.3) {4};
    \draw[graphedge] (v1)--node[above left,font=\small,fill=white,inner sep=1pt]{#1}(v2);
    \draw[graphedge] (v1)--node[below left,font=\small,fill=white,inner sep=1pt]{#2}(v3);
    \draw[graphedge] (v3)--node[right,font=\small,fill=white,inner sep=1pt]{#3}(v2);
    \draw[graphedge] (v2)--node[above right,font=\small,fill=white,inner sep=1pt]{#4}(v4);
    \draw[graphedge] (v3)--node[below right,font=\small,fill=white,inner sep=1pt]{#5}(v4);
  \end{tikzpicture}}

\title{{\vnheavy L4c}: Shortest-Path Algorithms}
\subtitle{CHEME 5800 Cornell University}
\author{Copyright Jeffrey Varner 2026. All Rights Reserved}

\begin{document}

% ==== title ==================================================
\vntitlepage

% ==== objectives =============================================
\begin{frame}{Objectives for Today}
  Today we find the least-cost routes from one starting vertex and compare
  how Dijkstra and Bellman-Ford calculate them.

  \vspace{0.9em}
  {\large\textbf{\ink{Today's concepts}}}
  \vspace{0.4em}
  \begin{itemize}\setlength{\itemsep}{0.7em}
    \item \hilite{Shortest paths and relaxation}: calculate path cost and update
      distance and predecessor maps.
    \item \hilite{Dijkstra's algorithm}: explain why nonnegative edge weights
      allow a selected distance to be declared final.
    \item \hilite{Bellman-Ford}: explain repeated relaxation, negative edges,
      and detection of negative cycles reachable from the source.
  \end{itemize}

  \slidenote{Lecture:}{\href{https://github.com/varnerlab/CHEME-5800-CourseRepository-Fall-2026/blob/main/weeks/week-04/L4c/CHEME-5800-L4c-Lecture-ShortestPathAlgorithms-Fall-2026.ipynb}{Shortest-Path Algorithms}}
\end{frame}

\begin{frame}{Lecture and Lab}
  \textbf{\ink{Lecture:}}
  \href{https://github.com/varnerlab/CHEME-5800-CourseRepository-Fall-2026/blob/main/weeks/week-04/L4c/CHEME-5800-L4c-Lecture-ShortestPathAlgorithms-Fall-2026.ipynb}{Shortest-Path Algorithms}

  \vspace{0.8em}
  Define route cost, follow both algorithms step by step, and distinguish a
  negative edge that lowers cost from a negative cycle that can be repeated
  to keep lowering it.

  \vspace{1.1em}
  \textbf{\ink{Lab L4d:}}
  \href{https://github.com/varnerlab/CHEME-5800-CourseRepository-Fall-2026/blob/main/weeks/week-04/L4d/CHEME-5800-L4d-Lab-ProductionPlanningShortestPath-Fall-2026.ipynb}{Production Planning with Shortest Paths}

  \vspace{0.8em}
  Predict a least-cost production route, compute it with Dijkstra, verify it
  with Bellman-Ford, and determine how large a discount must be to change the
  least-cost route.
\end{frame}

% ==== problem formulation ====================================
\begin{frame}{Which Route Has the Lowest Total Cost?}
  An edge weight can represent travel time, communication delay, monetary cost,
  or the number of unit operations in a chemical process.

  \vspace{0.9em}
  The common question is:

  \vspace{0.5em}
  \begin{center}
    \begin{minipage}{0.82\textwidth}
      \centering\Large
      Among all feasible routes from a source to a target, which route has the
      \hilite{smallest total edge weight}?
    \end{minipage}
  \end{center}

  \vspace{0.9em}
  Checking every route separately is usually impractical: even a graph with
  no cycles can have exponentially many routes.
\end{frame}

\begin{frame}{A Feasible Route Follows Directed Edges}
  Let $\mathcal G=(\mathcal V,\mathcal E)$ be a finite directed graph with vertex
  set $\mathcal V$ and edge set $\mathcal E$. The function
  $w:\mathcal E\rightarrow\mathbb R$ assigns each edge a real-valued cost.

  We write $|\mathcal V|$ for the number of vertices and $|\mathcal E|$ for the
  number of edges. A route $P$ from source $s$ to target $t$ with $k$ edges has
  the ordered vertices:
  \[
    P=\langle v_0,v_1,\ldots,v_k\rangle,
    \qquad v_0=s,\quad v_k=t,
  \]
  with $(v_i,v_{i+1})\in\mathcal E$ for $i=0,\ldots,k-1$.

  Its total weight is the sum of its edge weights:
  \[
    \underbrace{w(P)=\sum_{i=0}^{k-1}w(v_i,v_{i+1})}_{\text{route cost}}.
  \]
\end{frame}

\begin{frame}{Shortest-Path Distance Measures Minimum Cost}
  Let $\mathcal P_{s,t}$ be the set of feasible routes from $s$ to $t$.
  A negative-weight cycle is a directed cycle whose edge weights sum to a
  negative number. If no such cycle lies on an $s$-to-$t$ route, define:
  \[
    d(s,t)=
    \begin{cases}
      \displaystyle\min_{P\in\mathcal P_{s,t}}w(P),
        & \mathcal P_{s,t}\neq\varnothing,\\[5pt]
      +\infty, & \mathcal P_{s,t}=\varnothing.
    \end{cases}
  \]

  An unreachable target has $d(s,t)=+\infty$. For a finite minimum,
  $\min$ gives the cost and $\arg\min$ gives the routes with that cost.
  Choose one such route and call it $P^*$:
  \[
    P^*\in\arg\min_{P\in\mathcal P_{s,t}}w(P),
    \qquad w(P^*)=d(s,t).
  \]
\end{frame}

\begin{frame}{Record Distances and Predecessors}
  After finding shortest paths from source $s$, we record two values for each
  vertex $v\in\mathcal V$:
  \[
    \underbrace{\operatorname{dist}[v]=d(s,v)}_{\text{minimum route cost}},
    \qquad
    \underbrace{\operatorname{prev}[v]}_{\text{previous vertex on the chosen route}}.
  \]

  \vspace{0.7em}
  \begin{itemize}\setlength{\itemsep}{0.75em}
    \item \texttt{nothing} means no predecessor is recorded; the source has
      $\operatorname{prev}[s]=\texttt{nothing}$.
    \item Follow predecessors backward from a target, then reverse the sequence,
      to reconstruct one minimum-cost route.
    \item If several routes tie, the predecessor map may record any one of them.
    \item An unreachable vertex keeps distance $+\infty$ and no predecessor.
  \end{itemize}
\end{frame}

\begin{frame}{A Shortest Path Contains Shortest Paths}
  Suppose a shortest path from source $s$ to target $t$ passes through vertex $u$.

  \vspace{0.8em}
  The part of this path from $s$ to $u$ must also be a shortest path from $s$ to $u$.

  \vspace{0.8em}
  If we could reach $u$ more cheaply, we could take that cheaper route and then
  follow the original path from $u$ to $t$. This would lower the total cost to
  $t$, contradicting our assumption that the original path was shortest.

  \vspace{0.9em}
  \begin{center}
    \Large
    This property is called \hilite{optimal substructure}.
  \end{center}

  \slidenote{Why this helps:}{extend the cheapest route found to $u$ to look for cheaper routes to its neighbors.}
\end{frame}

% ==== relaxation =============================================
\begin{frame}{Relaxation Starts with One Known Route}
  At initialization, the zero-edge route from $s$ to itself is the only known route:
  \[
    \operatorname{dist}[v]=
    \begin{cases}
      0, & v=s,\\
      +\infty, & v\neq s,
    \end{cases}
    \qquad
    \operatorname{prev}[v]=\texttt{nothing}.
  \]

  \vspace{0.8em}
  A finite estimate $\operatorname{dist}[v]$ is always the weight of an actual
  source-to-$v$ route discovered so far.

  \vspace{0.6em}
  During the calculation, $\operatorname{dist}[v]$ is the lowest cost found so far.
  It can decrease when we find a cheaper route.
\end{frame}

\begin{frame}{Relaxation Checks for a Cheaper Route}
  Suppose we know a route to $u$. Following edge $(u,v)\in\mathcal E$
  extends that route to $v$.
  We use $\operatorname{alt}(u,v)$, short for \hilite{alternative route cost},
  for the cost of reaching $v$ this way:
  \[
    \underbrace{\operatorname{alt}(u,v)=
      \operatorname{dist}[u]+w(u,v)}_{\text{candidate route cost to }v}.
  \]

  Relaxation updates both maps only when the candidate is strictly cheaper:
  \[
  (\operatorname{dist}[v],\operatorname{prev}[v])\leftarrow
  \begin{cases}
    (\operatorname{alt}(u,v),u),
      & \operatorname{alt}(u,v)<\operatorname{dist}[v],\\[3pt]
    (\operatorname{dist}[v],\operatorname{prev}[v]),
      & \operatorname{alt}(u,v)\geq\operatorname{dist}[v].
  \end{cases}
  \]

  In the algorithm, $\operatorname{alt}$ stores this cost for the current edge.
  Tied costs leave the predecessor unchanged.
\end{frame}

\begin{frame}{Example: A Four-Vertex Graph}
  \begin{center}
  \begin{tikzpicture}[x=1.65cm,y=1.15cm]
    \node[graphsource] (v1) at (0,1.3) {1};
    \node[graphvertex] (v3) at (2,0) {3};
    \node[graphvertex] (v2) at (2,2.6) {2};
    \node[graphvertex] (v4) at (4.3,1.3) {4};
    \draw[graphedge] (v1)--node[above left,font=\small,fill=white,inner sep=1pt]{4}(v2);
    \draw[graphoptimal] (v1)--node[below left,font=\small,fill=white,inner sep=1pt]{1}(v3);
    \draw[graphoptimal] (v3)--node[right,font=\small,fill=white,inner sep=1pt]{2}(v2);
    \draw[graphoptimal] (v2)--node[above right,font=\small,fill=white,inner sep=1pt]{1}(v4);
    \draw[graphedge] (v3)--node[below right,font=\small,fill=white,inner sep=1pt]{5}(v4);
  \end{tikzpicture}
  \end{center}

  Edge labels give costs. Red edges mark the minimum-cost route from source 1
  to target 4. We first examine how relaxation improves the estimate for vertex 2.
\end{frame}

\begin{frame}{Example: Relaxing One Edge}
  From source 1, the direct edges give $\operatorname{dist}[2]=4$ and
  $\operatorname{dist}[3]=1$. Both vertices initially have predecessor 1.

  \vspace{0.7em}
  Can we reach vertex 2 more cheaply through vertex 3? For edge $(3,2)$,
  the alternative route cost is:
  \[
    \operatorname{alt}(3,2)=\operatorname{dist}[3]+w(3,2)=1+2=\hilite{3}.
  \]

  Since $3<4$, we update both maps:
  \[
    \operatorname{dist}[2]\leftarrow3,\qquad
    \operatorname{prev}[2]\leftarrow3.
  \]

  Following predecessors from vertex 2 leads to vertex 3 and then vertex 1.
  Reversing this sequence gives the cheaper route $1\rightarrow3\rightarrow2$.
  We still need to establish when a distance estimate is final.
\end{frame}

\begin{frame}{What Does Relaxation Guarantee?}
  Any route we discover costs at least as much as the shortest route.
  For a vertex $v$ with a finite shortest-path distance:
  \[
    \underbrace{d(s,v)\leq\operatorname{dist}[v]}_{\text{estimate is at least the minimum}}.
  \]

  Immediately after relaxing $(u,v)$, with $u\neq v$, the estimate for $v$
  is no greater than the cost of reaching it through $u$:
  \[
    \underbrace{\operatorname{dist}[v]
      \leq\operatorname{dist}[u]+w(u,v)}_{\text{after checking the route through }u}.
  \]

  \vspace{0.7em}
  If $\operatorname{dist}[u]$ later decreases, we may need to check $(u,v)$ again.
\end{frame}

\transitionframe{Dijkstra and Bellman-Ford use the \hilite{same relaxation rule}.
Dijkstra selects a vertex and checks its outgoing edges.
Bellman-Ford checks every edge repeatedly.}

% ==== Dijkstra ===============================================
\begin{frame}{Dijkstra Requires Nonnegative Edge Weights}
  Dijkstra selects the smallest distance estimate among vertices not yet selected.
  For that estimate to be final, every edge must satisfy:
  \[
    w(u,v)\geq0\qquad\text{for every }(u,v)\in\mathcal E.
  \]

  Initialize:
  \begin{itemize}\setlength{\itemsep}{0.55em}
    \item the distance and predecessor maps defined earlier;
    \item an empty set $\mathcal S$ of selected vertices;
    \item a min-priority queue $\mathcal Q$, using distance estimates as priorities.
  \end{itemize}

  Start with only $s$ in $\mathcal Q$, at priority zero.
  The queue returns the vertex with the smallest estimate first.
\end{frame}

\begin{frame}{Dijkstra Finds One Shortest Distance at a Time}
  While the min-priority queue $\mathcal Q$ is not empty:

  \vspace{0.6em}
  \begin{enumerate}\setlength{\itemsep}{0.65em}
    \item Remove the vertex $u$ with the smallest distance estimate from $\mathcal Q$.
    \item If $u\in\mathcal S$, skip it; otherwise add $u$ to $\mathcal S$.
      Its distance is now final.
    \item For every outgoing edge $(u,v)$, compute
      $\operatorname{alt}=\operatorname{dist}[u]+w(u,v)$.
    \item If $\operatorname{alt}<\operatorname{dist}[v]$, update the distance and
      set $\operatorname{prev}[v]=u$. Insert $v$ into $\mathcal Q$, or lower its
      priority to $\operatorname{alt}$ if it is already queued.
  \end{enumerate}

  The adjacency list gives the outgoing edges and their weights.
\end{frame}

\begin{frame}{Dijkstra Trace on the Worked Graph}
  Each row follows relaxation from the selected vertex.
  \hilite{Coral} marks changed estimates; the source distance stays zero.

  \vspace{0.8em}
  \begin{columns}[c,onlytextwidth]
    \begin{column}{0.29\textwidth}
      \centering
      \tracegraph{4}{1}{2}{1}{5}

      \small Source 1; edge labels give costs.
    \end{column}
    \begin{column}{0.68\textwidth}
      \small
      \setlength{\tabcolsep}{2.5pt}
      \renewcommand{\arraystretch}{1.3}
      \begin{tabularx}{\linewidth}{@{}lcccY@{}}
        \toprule
        Selected & $\operatorname{dist}[2]$ & $\operatorname{dist}[3]$ & $\operatorname{dist}[4]$ & Queue \\
        \midrule
        Initial & $\infty$ & $\infty$ & $\infty$ & \texttt{1:0} \\
        1 & \hilite{4} & \hilite{1} & $\infty$ & \texttt{3:1, 2:4} \\
        3 & \hilite{3} & 1 & \hilite{6} & \texttt{2:3, 4:6} \\
        2 & 3 & 1 & \hilite{4} & \texttt{4:4} \\
        4 & 3 & 1 & 4 & Empty \\
        \bottomrule
      \end{tabularx}

      \vspace{0.5em}
      Queue entries are \texttt{vertex:priority}, in increasing priority order.
    \end{column}
  \end{columns}

  \vspace{0.8em}
  The predecessors are $\operatorname{prev}[4]=2$,
  $\operatorname{prev}[2]=3$, and $\operatorname{prev}[3]=1$.
  Reverse the backward sequence to recover $1\rightarrow3\rightarrow2\rightarrow4$.
\end{frame}

\begin{frame}{Why Can't a Selected Distance Improve Later?}
  The source's shortest-path distance is zero. Assume all previously selected
  distances are correct. Suppose a route from $s$ to the next vertex $u$
  costs less than $\operatorname{dist}[u]$.

  \begin{center}
  \begin{tikzpicture}[x=1.3cm,y=0.8cm]
    \node[graphsource,fill=vnaccent!10] (s) at (0,0) {$s$};
    \node[graphsource,fill=vnaccent!10] (x) at (2,0) {$x$};
    \node[graphvertex] (y) at (4,0) {$y$};
    \node[graphvertex] (u) at (6,0) {$u$};
    \draw[graphedge,dashed] (s)--(x);
    \draw[graphoptimal] (x)--(y);
    \draw[graphedge,dashed] (y)--(u);
    \draw[vnrule,dashed] (3,-0.65)--(3,1.15);
    \node[above,font=\small,fill=white,inner sep=1pt] at (3,0.1) {relaxed};
    \node[font=\small\vnmedium,text=vnslidebody] at (1,1.0) {Already selected};
    \node[font=\small\vnmedium,text=vnslidebody] at (5,1.0) {Not yet selected};
    \node[font=\small,text=vnslidebody] at (3,-1.0) {Dashed arrows may represent several edges.};
  \end{tikzpicture}
  \end{center}

  Let $y$ be the first vertex on this route not yet selected, and $x$ its predecessor.
  Relaxing $(x,y)$ gave $y$ an estimate no greater than the cost of reaching $y$
  along this route. The remaining edges add no negative cost, so:
  \[
    \operatorname{dist}[y]
    \leq\text{cost of the supposed cheaper route}
    <\operatorname{dist}[u].
  \]

  This contradicts our choice of $u$.
  Hence $\operatorname{dist}[u]=d(s,u)$.
\end{frame}

\begin{frame}{Negative Edges Can Change a Selected Distance}
  Consider a different graph with source vertex 1 and three directed edges:
  \[
    w(1,2)=2,\qquad w(1,3)=5,\qquad w(3,2)=-4.
  \]

  \vspace{0.8em}
  If we apply Dijkstra's selection rule despite the negative edge, vertex 2 is
  selected with distance 2 before vertex 3.

  \vspace{0.8em}
  Yet the route through vertex 3 has lower cost:
  \[
    1\rightarrow3\rightarrow2,\qquad 5-4=\hilite{1}.
  \]

  Thus, a selected distance need not be final. Our implementation rejects
  negative-edge inputs before applying the selection rule.
\end{frame}

\begin{frame}{Dijkstra: Computational Cost}
  With a weighted adjacency list and a binary-heap min-priority queue,
  Dijkstra runs in:
  \[
    \underbrace{\mathcal O\!\left((|\mathcal V|+|\mathcal E|)
      \log|\mathcal V|\right)}_{\text{Dijkstra running time}}.
  \]

  \vspace{0.8em}
  \begin{itemize}\setlength{\itemsep}{0.75em}
    \item Each reachable vertex is selected once.
    \item Each outgoing edge is examined when its starting vertex is selected.
    \item Queue insertions, removals, and priority updates each take
      $\mathcal O(\log|\mathcal V|)$ time.
  \end{itemize}

  \slidenote{Scope:}{the bound matches the course implementation using adjacency lists and a binary-heap priority queue.}
\end{frame}

% ==== Bellman-Ford ===========================================
\begin{frame}{Bellman-Ford Repeatedly Relaxes Every Edge}
  Bellman-Ford checks every edge during each \hilite{pass}. Repeating these checks
  allows later discoveries to lower earlier distance estimates, including when
  some weights are negative.

  \vspace{0.7em}
  Initialize the same distance and predecessor maps. Then repeat at most
  $|\mathcal V|-1$ times:
  \begin{enumerate}\setlength{\itemsep}{0.55em}
    \item Visit every edge $(u,v)\in\mathcal E$.
    \item If $\operatorname{dist}[u]$ is finite and the candidate is cheaper,
      relax $(u,v)$.
    \item If no distance changes during a complete pass, stop.
  \end{enumerate}

  Each update is available to edges checked later in the same pass.
  Thus, edge order can change how many passes are needed.
\end{frame}

\begin{frame}{Bellman-Ford: Following One Pass}
  Starting at source 1, visit the edges in the listed order.
  \hilite{Coral} marks an estimate changed by that relaxation.

  \vspace{0.7em}
  \begin{columns}[c,onlytextwidth]
    \begin{column}{0.29\textwidth}
      \centering
      \tracegraph{4}{1}{2}{1}{5}

      \small $\operatorname{dist}[1]=0$ throughout.
    \end{column}
    \begin{column}{0.66\textwidth}
      \small
      \setlength{\tabcolsep}{5pt}
      \renewcommand{\arraystretch}{1.25}
      \begin{tabularx}{\linewidth}{@{}Yccc@{}}
        \toprule
        Edge relaxed & $\operatorname{dist}[2]$ & $\operatorname{dist}[3]$ & $\operatorname{dist}[4]$ \\
        \midrule
        Initial & $\infty$ & $\infty$ & $\infty$ \\
        $(1,2)$ & \hilite{4} & $\infty$ & $\infty$ \\
        $(1,3)$ & 4 & \hilite{1} & $\infty$ \\
        $(3,2)$ & \hilite{3} & 1 & $\infty$ \\
        $(2,4)$ & 3 & 1 & \hilite{4} \\
        $(3,4)$ & 3 & 1 & 4 \\
        \bottomrule
      \end{tabularx}
    \end{column}
  \end{columns}

  \vspace{0.7em}
  Relaxing $(3,2)$ lowers the estimate for vertex 2 from 4 to 3.
  The next edge, $(2,4)$, uses that improvement immediately to give $3+1=4$.
\end{frame}

\begin{frame}{Bellman-Ford: Comparing Edge Orders}
  Reverse the preceding edge order:
  \[
    (3,4),\quad(2,4),\quad(3,2),\quad(1,3),\quad(1,2).
  \]

  Each triple gives $(\operatorname{dist}[2],\operatorname{dist}[3],\operatorname{dist}[4])$.
  \hilite{Coral} marks changes during that pass.

  \vspace{0.6em}
  \begin{columns}[c,onlytextwidth]
    \begin{column}{0.29\textwidth}
      \centering
      \tracegraph{4}{1}{2}{1}{5}
    \end{column}
    \begin{column}{0.66\textwidth}
      \small
      \renewcommand{\arraystretch}{1.35}
      \begin{tabularx}{\linewidth}{@{}l@{\hspace{1em}}Y@{\hspace{1em}}Y@{}}
        \toprule
        Pass & Listed order & Reversed order \\
        \midrule
        0 & $(\infty,\infty,\infty)$ & $(\infty,\infty,\infty)$ \\
        1 & $(\hilite{3},\hilite{1},\hilite{4})$ & $(\hilite{4},\hilite{1},\infty)$ \\
        2 & $(3,1,4)$; stop & $(\hilite{3},1,\hilite{5})$ \\
        3 & Stopped & $(3,1,\hilite{4})$ \\
        \bottomrule
      \end{tabularx}
    \end{column}
  \end{columns}

  \vspace{0.7em}
  The listed order stops after pass 2 makes no changes.
  Reversing the order puts $(2,4)$ before $(3,2)$, so vertex 4 improves one pass later.
\end{frame}

\begin{frame}{Why $|\mathcal V|-1$ Passes Are Enough}
  Assume no negative-weight cycle is reachable from the source. After $k$ passes,
  distances are correct for every vertex with a shortest path containing at most
  $k$ edges.

  \vspace{0.8em}
  If a route from the source repeats a vertex, the intervening edges form a cycle.
  Removing that cycle cannot increase the cost because its total weight is nonnegative.

  \vspace{0.8em}
  We can therefore choose a shortest path with no repeated vertices, called a
  \hilite{simple path}. It visits at most $|\mathcal V|$ vertices and contains at most:
  \[
    \underbrace{|\mathcal V|-1}_{\text{maximum edges in a simple path}}
  \]

  This is why $|\mathcal V|-1$ relaxation passes are sufficient.
\end{frame}

\begin{frame}{Detecting a Reachable Negative Cycle}
  After the relaxation passes, check every edge once more.

  \vspace{0.7em}
  If an edge from a reachable vertex can still lower a distance estimate, there
  is a reachable cycle $C$ with total weight $w(C)<0$. For a route $P$ through $C$,
  repeating the cycle gives costs:
  \[
    w(P),\quad w(P)+w(C),\quad w(P)+2w(C),\quad\ldots
    \longrightarrow -\infty
    \qquad\text{when }w(C)<0.
  \]

  For any target reachable from that cycle, repeating the cycle keeps making
  the route cheaper. There is no minimum cost.

  \vspace{0.7em}
  A negative cycle unreachable from $s$ is skipped because its vertices keep
  distance $+\infty$.
\end{frame}

\begin{frame}{Bellman-Ford: Computational Cost}
  In the worst case, Bellman-Ford scans every edge during each of
  $|\mathcal V|-1$ passes and during the cycle check:
  \[
    \underbrace{\mathcal O(|\mathcal V|\,|\mathcal E|)}_{\text{running time}},
    \qquad
    \underbrace{\mathcal O(|\mathcal V|)}_{\text{distance and predecessor storage}}.
  \]

  \vspace{0.8em}
  We can stop before the maximum number of passes if a complete pass changes no distances.

  \vspace{0.8em}
  Bellman-Ford can find shortest paths with negative edges and report when a
  reachable negative cycle prevents a minimum cost.
\end{frame}

\begin{frame}{Choosing Between Dijkstra and Bellman-Ford}
  \renewcommand{\arraystretch}{1.18}
  \begin{tabularx}{\textwidth}{@{}>{\vnmedium}p{0.26\textwidth}@{\hspace{1em}}Y@{\hspace{1em}}Y@{}}
    \toprule
    Question & Dijkstra & Bellman-Ford \\
    \midrule
    Weights & All nonnegative & May be negative \\
    \addlinespace[0.5em]
    Update rule & Select smallest estimate; relax outgoing edges & Relax every edge repeatedly \\
    \addlinespace[0.5em]
    Negative cycles & Not supported & Detects reachable negative cycles \\
    \addlinespace[0.5em]
    Worst-case time & $\mathcal O((|\mathcal V|+|\mathcal E|)\log|\mathcal V|)$ & $\mathcal O(|\mathcal V||\mathcal E|)$ \\
    \addlinespace[0.5em]
    Storage used here & Weighted adjacency list & Edge list \\
    \bottomrule
  \end{tabularx}

  \vspace{0.7em}
  Our Bellman-Ford implementation reports an error for a reachable negative cycle.
\end{frame}

% ==== worked comparison ======================================
\begin{frame}{Both Algorithms Find the Same Minimum Costs}
  On our four-vertex graph, both algorithms give the same shortest-path distances:
  \[
    \operatorname{dist}_{\mathrm{Dijkstra}}[v]
    =\operatorname{dist}_{\mathrm{Bellman-Ford}}[v]
    \qquad\text{for every vertex }v.
  \]

  Both predecessor maps reconstruct the unique minimum-cost route to vertex 4:
  \[
    P^*=\langle1,3,2,4\rangle,
    \qquad w(P^*)=1+2+1=4.
  \]

  \vspace{0.8em}
  All weights are nonnegative, so both algorithms apply. They use different
  orders and numbers of edge checks to find the same minimum costs.
\end{frame}

\begin{frame}{Breadth-First Search When Every Edge Costs 1}
  If every edge costs 1, a route's cost equals its number of edges.
  Breadth-first search, used in L4b, finds routes with the fewest edges.

  \vspace{0.8em}
  \begin{columns}[c,onlytextwidth]
    \begin{column}{0.29\textwidth}
      \centering
      \tracegraph{1}{1}{1}{1}{1}

      \small Same connections; every weight is now 1.
    \end{column}
    \begin{column}{0.66\textwidth}
      \renewcommand{\arraystretch}{1.4}
      \begin{tabularx}{\linewidth}{@{}Y@{\hspace{1em}}r@{}}
        \toprule
        Route from 1 to 4 & Cost \\
        \midrule
        $1\rightarrow2\rightarrow4$ & \hilite{2} \\
        $1\rightarrow3\rightarrow4$ & \hilite{2} \\
        $1\rightarrow3\rightarrow2\rightarrow4$ & 3 \\
        \bottomrule
      \end{tabularx}
    \end{column}
  \end{columns}

  \vspace{0.9em}
  The two-edge routes now tie for the minimum cost of 2.
  With the original weights, the three-edge route cost 4 and was cheapest.
\end{frame}

\begin{frame}{A Shortest Path Can Include a Negative Edge}
  \begin{center}
  \begin{tikzpicture}[x=1.7cm,y=1.0cm]
    \node[graphsource] (v1) at (0,1.2) {1};
    \node[graphvertex] (v2) at (1.8,2.2) {2};
    \node[graphvertex] (v3) at (3.6,1.2) {3};
    \node[graphvertex] (v4) at (5.4,1.2) {4};
    \draw[graphoptimal] (v1)--node[above left,font=\small,fill=white,inner sep=1pt]{4}(v2);
    \draw[graphedge] (v1)--node[below,font=\small,fill=white,inner sep=1pt]{5}(v3);
    \draw[graphnegative] (v2)--node[above right,font=\small,fill=white,inner sep=1pt]{$-2$}(v3);
    \draw[graphoptimal] (v3)--node[above,font=\small,fill=white,inner sep=1pt]{3}(v4);
  \end{tikzpicture}
  \end{center}

  Bellman-Ford finds the minimum-cost route:
  \[
    1\rightarrow2\rightarrow3\rightarrow4,
    \qquad 4-2+3=\hilite{5}.
  \]

  This graph has no cycles, so the negative edge cannot be repeated to keep
  lowering the cost.
\end{frame}

\begin{frame}{A Negative Cycle Keeps Lowering Cost}
  \begin{center}
  \begin{tikzpicture}[x=1.6cm,y=1.25cm]
    \node[graphsource] (v1) at (0,0) {1};
    \node[graphvertex] (v2) at (2.4,1.7) {2};
    \node[graphvertex] (v3) at (4.8,0) {3};
    \draw[graphoptimal] (v1)--node[above left,font=\small,fill=white,inner sep=1pt]{1}(v2);
    \draw[graphnegative] (v2)--node[above right,font=\small,fill=white,inner sep=1pt]{$-2$}(v3);
    \draw[graphoptimal] (v3)--node[below,font=\small,fill=white,inner sep=1pt]{0}(v1);
  \end{tikzpicture}
  \end{center}

  The cycle weight is:
  \[
    w(C)=1-2+0=\hilite{-1}.
  \]
  Each trip around the cycle lowers the route cost by one. Bellman-Ford's extra
  check finds another possible improvement and reports an error: there is no
  minimum-cost route to any vertex in this graph.
\end{frame}

\begin{frame}{Check the Routes, Costs, and Expected Errors}
  The notebook checks the calculated results and the expected errors:

  \vspace{0.7em}
  \begin{itemize}\setlength{\itemsep}{0.75em}
    \item \textbf{\ink{Nonnegative graph:}} both algorithms return
      $\langle1,3,2,4\rangle$ with target cost 4 and identical distance maps.
    \item \textbf{\ink{Negative edge, no negative cycle:}} Bellman-Ford returns
      $\langle1,2,3,4\rangle$ with target cost 5.
    \item \textbf{\ink{Expected errors:}} Dijkstra raises an error for the
      negative edge; Bellman-Ford raises an error for the reachable negative cycle.
  \end{itemize}

  \slidenote{What we check:}{all distances on the nonnegative graph, selected routes and costs, and expected errors.}
\end{frame}

% ==== lab bridge =============================================
\begin{frame}{Lab L4d Applies Both Algorithms}
  The lab represents production steps and their costs as a weighted directed graph.

  \vspace{0.7em}
  \begin{enumerate}\setlength{\itemsep}{0.7em}
    \item Predict the least-cost route by inspecting the process costs.
    \item Compute the route with Dijkstra because all costs are nonnegative.
    \item Verify the distance and route with Bellman-Ford.
    \item Discount one production step and find the step cost at which the two
      routes have equal total cost.
  \end{enumerate}

  \vspace{0.6em}
  We then examine which route is cheaper above and below that step cost.
\end{frame}

% ==== summary ================================================
\begin{frame}{Summary}
  We compared two algorithms for finding minimum-cost routes.

  \vspace{0.7em}
  {\large\textbf{\ink{Key takeaways}}}
  \vspace{0.35em}
  \begin{itemize}\setlength{\itemsep}{0.6em}
    \item \textbf{\ink{Route costs and relaxation:}} We summed edge weights,
      improved distance estimates by relaxation, and used predecessors to reconstruct routes.
    \item \textbf{\ink{Dijkstra's algorithm:}} We used a priority queue to select
      the smallest estimate and showed why nonnegative weights make it final.
    \item \textbf{\ink{Bellman-Ford's algorithm:}} We used repeated relaxation
      to handle negative edges and an extra scan to detect reachable negative cycles.
  \end{itemize}

  \vspace{0.5em}
  In \href{https://github.com/varnerlab/CHEME-5800-CourseRepository-Fall-2026/blob/main/weeks/week-04/L4d/CHEME-5800-L4d-Lab-ProductionPlanningShortestPath-Fall-2026.ipynb}{the L4d lab},
  we apply these methods to production costs and equipment discounts.
\end{frame}

\transitionframe{That's a wrap. \hilite{Which property of the edge weights tells
you which shortest-path algorithm is valid?}}

\end{document}
